\subsection{PROJECTIVE LINEAR VARIETIES}

Let W be a vector subspace of a vector space V; the canonical image of \(W - \{0\}\) in the projective space \(\mathbf{P}(V)\) derived from V is called a projective linear variety (or simply a linear variety when no confusion is to be feared); as the equivalence relation \(\Delta(W)\) on \(W - \{0\}\) is induced by the relation \(\Delta(V)\) , the projective linear variety the image of \(W - \{0\}\) in \(\mathbf{P}(V)\) can be identified with the projective space \(\mathbf{P}(W)\) derived from W and hence we may speak of the dimension of such a variety. In a projective space \(\mathbf{P}(V)\) , the canonical image of a hyperplane (with the origin removed) of V is a linear variety called a projective hyperplane (or simply a hyperplane); if \(\mathbf{P}(V)\) is of finite dimension n the hyperplanes in \(\mathbf{P}(V)\) are the linear varieties of dimension n - 1.

Every proposition concerning vector subspaces of a vector space goes over to a proposition concerning projective linear varieties. For example, if a projective space \(\mathbf{P}(\mathbf{V})\) is of finite dimension n and \((e_{i})_{0\leqslant i\leqslant n}\) is a basis of V, every linear variety \(\mathbf{L}\subset\mathbf{P}(\mathbf{V})\) of dimension r can be defined by a system of n-r homogeneous linear equations

\[
\sum_ {i = 0} ^ {n} \xi_ {i} \alpha_ {i j} = 0 \quad (1 \leqslant j \leqslant n - r)\tag{4}
\]

between the homogeneous coordinates \(\xi_{i}\) ( \(0 \leqslant i \leqslant n\) ) of a point of \(\mathbf{P}(\mathbf{V})\) with respect to the basis \((e_{i})\) , the left hand sides of (4) being linearly independent forms on V. In particular, a projective hyperplane is defined by a single homogeneous linear equation with coefficients not all zero. Conversely, the points of \(\mathbf{P}(\mathbf{V})\) satisfying an arbitrary system of homogeneous linear equations with respect to the \(\xi_{i}\) form a linear variety L; if the system considered consists of \(k \leqslant n + 1\) equations, L is of dimension \(\geqslant n - k\) .

Every intersection of linear varieties of \(\mathbf{P}(\mathbf{V})\) is a linear variety; for every subset A of \(\mathbf{P}(\mathbf{V})\) , there exists a smallest linear variety L containing A; it is called the linear variety generated by A and A is called a generating system of L.

If W is the vector subspace of V generated by \(\pi^{1}(A)\) , then L = P(W).

If \(\mathbf{L}\) and \(\mathbf{M}\) are any two linear varieties in \(\mathbf{P}(\mathbf{V})\) and \(\mathbf{N}\) the variety generated by \(\mathbf{L} \cup \mathbf{M}\) , then (§ 7, no. 3, Corollary 3 to Proposition 4)

\[
\dim \mathbf {L} + \dim \mathbf {M} = \dim (\mathbf {L} \cap \mathbf {M}) + \dim \mathbf {N}.\tag{5}
\]

In particular, if \(\mathbf{P}(\mathbf{V})\) is finite-dimensional and \(\dim \mathbf{L} + \dim \mathbf{M} \geqslant \dim \mathbf{P}(\mathbf{V})\) , it follows from (5) that \(\mathbf{L} \cap \mathbf{M}\) is non-empty.

Let \((x_{i}), (y_{i})\) be two families of points in the vector space V with the same indexing set, such that \(y_{i} = \lambda_{i}x_{i}\) , where \(\lambda_{i}^{*} \neq 0\) for all i. If the family \((x_{i})\) is free, so is \((y_{i})\) and conversely; then it is said that the family of points \(\pi(x_{i})\) of P(V) is projectively free (or simply free). It amounts to the same to say that for every index K, the point \(\pi(x_{K})\) does not belong to the linear variety generated by the \(\pi(x_{i})\) for \(i \neq K\) . A family of points of P(V) which is not projectively free is called projectively related (or simply related).

For a family \((x_{i})\) of points of \(\mathbf{V} - \{0\}\) to be such that the family \((\pi (x_{i}))\) is projectively free and generates \(\mathbf{P}(\mathbf{V})\) , it is necessary and sufficient that \((x_{i})\) be a basis of \(\mathbf{V}\) . If \(\mathbf{P}(\mathbf{V})\) is of dimension \(n\) the number of elements in such a family is therefore \(n + 1\) . Note that giving such a family \((\pi (x_{i}))\) in \(\mathbf{P}(\mathbf{V})\) does not determine (even to within a left factor) the homogeneous coordinates of a given point of \(\mathbf{P}(\mathbf{V})\) with respect to a basis \((y_{i})\) of \(\mathbf{V}\) such that \(\pi (y_{i}) = \pi (x_{i})\) for all \(\iota\) (cf.~no. 6).

